% 杨辉三角（帕斯卡三角）7 层
\documentclass[tikz,border=4pt]{standalone}
\usepackage{ctex}
\usepackage{amsmath}
\usetikzlibrary{calc}
\begin{document}
\begin{tikzpicture}[
    node distance=0pt,
    every node/.style={
      circle, draw=blue!50!black, fill=blue!10,
      minimum size=0.85cm, inner sep=0pt,
      font=\small
    },
    nlabel/.style={
      draw=none, fill=none, font=\footnotesize, gray!40!black,
      minimum size=0pt
    },
]
  % 行间距、列间距
  \def\dx{0.9}
  \def\dy{-0.95}

  % 数值（杨辉三角前 7 行）：第 n 行（0 起）的第 k 项 C(n,k)
  % 第 0 行：1
  % 第 1 行：1 1
  % 第 2 行：1 2 1
  % 第 3 行：1 3 3 1
  % 第 4 行：1 4 6 4 1
  % 第 5 行：1 5 10 10 5 1
  % 第 6 行：1 6 15 20 15 6 1
  \foreach \row/\vals in {
    0/{1},
    1/{1,1},
    2/{1,2,1},
    3/{1,3,3,1},
    4/{1,4,6,4,1},
    5/{1,5,10,10,5,1},
    6/{1,6,15,20,15,6,1}
  } {
    \foreach \v [count=\k from 0] in \vals {
      \pgfmathsetmacro{\xpos}{(\k - \row/2) * \dx}
      \pgfmathsetmacro{\ypos}{\row * \dy}
      \node at (\xpos, \ypos) {$\v$};
    }
  }

  % 行号标签（左侧）
  \foreach \row/\name in {0/{$n=0$}, 1/{$n=1$}, 2/{$n=2$}, 3/{$n=3$},
                          4/{$n=4$}, 5/{$n=5$}, 6/{$n=6$}} {
    \pgfmathsetmacro{\xpos}{(-\row/2 - 1.4) * \dx}
    \pgfmathsetmacro{\ypos}{\row * \dy}
    \node[nlabel] at (\xpos, \ypos) {\name};
  }

  % 标题（顶上）
  \node[nlabel, font=\small] at (0, 1) {\textbf{杨辉三角（帕斯卡三角）：}$C_n^k$ 即第 $n$ 行第 $k$ 个数};

  % 底部公式提示
  \node[nlabel, font=\small, align=center] at (0, 7.5*\dy)
    {每数 $=$ 上一行肩上两数之和：$C_n^k = C_{n-1}^{k-1} + C_{n-1}^{k}$};
\end{tikzpicture}
\end{document}
